\part{Smart pointers}
\frame{\partpage}

\begin{frame}{Motivation}
    \begin{itemize}
        \item Pointers in C++ are powerful, but with great power comes great responsibility! \pause
        \item In particular, responsibility for the \textbf{lifetime} of object instances,
            i.e.\ remembering to call \lstinline{delete} at the right time \pause
        \item C++11 introduced \textbf{smart pointers} to try and make this easier
    \end{itemize}
\end{frame}

\begin{frame}{Unique pointer}
    \begin{itemize}
        \item Defined in standard header file \lstinline{<memory>} \pause
        \item \lstinline{std::unique_ptr} has \textbf{sole ownership} of the instance to which it points \pause
        \item I.e.\ two \lstinline{unique_ptr}s \textbf{cannot} point to the same instance \pause
        \item Use \lstinline{std::make_unique<T>} instead of \lstinline{new T} \pause
        \item The instance is destroyed when the \lstinline{unique_ptr} goes \textbf{out of scope},
            or when it is assigned to point to a different instance
    \end{itemize}
\end{frame}

\begin{frame}[fragile]{Unique pointer example}
    \begin{lstlisting}
#include "stdafx.h" // --> #include <memory>

class MyClass { };

int main()
{
    std::unique_ptr<MyClass> p
        = std::make_unique<MyClass>();
    return 0;
    // p is automatically deleted here
}
    \end{lstlisting}
\end{frame}

\begin{frame}[fragile]{Aside: automatic type inference}
    \begin{itemize}
        \item When declaring a \textbf{local variable} with an initial value,
            can use \lstinline{auto} in place of the type \pause
        \item The compiler \textbf{infers} the type from the type of the initial value \pause
    \end{itemize}
    \begin{lstlisting}
// The following are equivalent:
std::unique_ptr<Duck> duck = std::make_unique<Duck>();
auto duck = std::make_unique<Duck>();
    \end{lstlisting} \pause
    \begin{itemize}
        \item Careful --- this can make your code \textbf{more readable} or \textbf{less readable}
            depending on how and where you use it!
    \end{itemize}
\end{frame}

\begin{frame}{Vectors of objects revisited}
    From last time: \pause
    \begin{itemize}
        \item \lstinline{std::vector<Duck>}: when the vector changes size, instances are 
            \textbf{copied} and the old ones \textbf{destroyed} \pause
        \item \lstinline{std::vector<Duck*>}: no copying and destroying, but you must
            remember to \lstinline{delete} elements on removal \pause
    \end{itemize}
    But now: \pause
    \begin{itemize}
        \item \lstinline{std::vector<std::unique_ptr<Duck>>}: class instances stay put,
            but are automatically destroyed when removed from the \lstinline{vector}!
    \end{itemize}
\end{frame}

\begin{frame}[fragile]{Unique pointers can't be copied}
    \begin{lstlisting}
auto p = std::make_unique<MyClass>();
auto q = p;
    \end{lstlisting} \pause
    \begin{itemize}
        \item This won't compile! \pause
    \end{itemize}
    \includegraphics[width=\textwidth]{unique_ptr_error} \pause
    \begin{itemize}
        \item Translation: assigning one \lstinline{unique_ptr} from another is forbidden
    \end{itemize}
\end{frame}

\begin{frame}[fragile]{Aside: default and deleted functions}
    \begin{itemize}
        \item As standard, C++ classes come with constructors and methods for \textbf{copying} and \textbf{moving} instances \pause
        \item You (the class developer) can \textbf{override} these methods, make them \textbf{protected} or \textbf{private},
            or \textbf{delete} them entirely (e.g.\ to forbid copying) \pause
        \item Further reading:
            \begin{itemize}
                \item \url{http://www.stroustrup.com/C++11FAQ.html#default}
                \item \url{http://en.cppreference.com/w/cpp/language/classes}
            \end{itemize}
    \end{itemize}
\end{frame}

\begin{frame}[fragile]{Unique pointers can be moved}
    \begin{lstlisting}
auto p = std::make_unique<MyClass>();
auto q = std::move(p);
    \end{lstlisting} \pause
    \begin{itemize}
        \item Now \lstinline{q} has \textbf{ownership} of the instance \pause
        \item \lstinline{p} no longer refers to this instance \pause
        \item In fact \lstinline{p} becomes null \pause
        \item Another way you are allowed to move \lstinline{unique_ptr}s is by \textbf{swapping}:
    \end{itemize}
    \begin{lstlisting}
auto p = std::make_unique<MyClass>("hello");
auto q = std::make_unique<MyClass>("world");
std::swap(p, q);
    \end{lstlisting}
\end{frame}

\begin{frame}{Shared pointers}
    \begin{itemize}
        \item Unique pointers have their uses, but uniqueness is a big restriction \pause
        \item C++ also has \lstinline{std::shared_ptr}, which does allow multiple pointers to the same
            instance \pause
        \item \textbf{Reference counting} ensures the instance is destroyed when
            \textbf{all \lstinline{shared_ptr}s} to it have gone away (and not before!)
    \end{itemize}
\end{frame}

\begin{frame}[fragile]
    \begin{lstlisting}
int main()
{
	bool flag = true;
	std::shared_ptr<MyClass> p = nullptr;
	
	if (flag)
	{
		auto q = std::make_shared<MyClass>();
		p = q;
	}
    /* q is out of scope, but p still holds a reference to the instance so it stays alive */
    
    if (p)
        p->doSomething();
    
    return 0;
    /* p is out of scope now, so the instance is destroyed */
}
    \end{lstlisting}
\end{frame}

\begin{frame}{Reference counting}
    \begin{itemize}
        \item Objects managed by \lstinline{shared_ptr} have a \textbf{reference count}, initially 0 \pause
        \item When a \lstinline{shared_ptr} is made to point to the instance,
            its reference count is \textbf{incremented} \pause
        \item When a \lstinline{shared_ptr} no longer points to the instance
            (because it goes out of scope or is reassigned),
            the reference count is \textbf{decremented} \pause
        \item When the reference count is decremented to 0,
            the instance is destroyed
    \end{itemize}
\end{frame}

\begin{frame}{Smart pointers and composition}
    \begin{itemize}
        \item We have seen \lstinline{unique_ptr} and \lstinline{shared_ptr} used as
            \textbf{local variables} \pause
        \item They can also be used as \textbf{fields} within classes \pause
        \item When the instance containing them is destroyed: \pause
            \begin{itemize}
                \item \lstinline{unique_ptr} destroys the instance to which it points \pause
                \item \lstinline{shared_ptr} decrements the reference count on the instance 
                    to which it points, which may result in it being destroyed \pause
            \end{itemize}
        \item ... just as you would expect (hopefully)!
    \end{itemize}
\end{frame}

\begin{frame}[fragile]
    \begin{lstlisting}
class Duck
{
public:
    Duck(const std::shared_ptr<Pond>& pond)
        : pond(pond),
          bill(std::make_unique<Bill>())
    { }
    
private:
    std::unique_ptr<Bill> bill;
    std::shared_ptr<Pond> pond;
};
    \end{lstlisting} \pause
    When a \lstinline{Duck} instance is destroyed: \pause
    \begin{itemize}
        \item Its \lstinline{Bill} is also destroyed; \pause
        \item Its \lstinline{Pond} is destroyed \textbf{if and only if} there are no other
            \lstinline{shared_ptr}s to it (e.g.\ in other \lstinline{Duck} instances)
    \end{itemize}
\end{frame}

\begin{frame}[fragile]
    \begin{lstlisting}
class Bar;

class Foo
{
public:
    std::shared_ptr<Bar> bar;
};

class Bar
{
public:
    std::shared_ptr<Foo> foo;
};

int main()
{
    auto p = std::make_shared<Foo>();
    p->bar = std::make_shared<Bar>();
    p->bar->foo = p;
}
    \end{lstlisting}
\end{frame}

\begin{frame}{Reference cycles}
    \begin{itemize}
        \item In the example on the previous slide,
            the instances are \textbf{not} destroyed properly at the end of \lstinline{main()} \pause
        \item \lstinline{Foo} still holds a reference to \lstinline{Bar},
            but \lstinline{Bar} holds a reference to \lstinline{Foo} \pause
        \item This prevents both objects' reference counts from reaching 0 \pause
        \item Solution: make one of the references a \textbf{weak pointer}
    \end{itemize}
\end{frame}

\begin{frame}[fragile]{Weak pointers}
    \begin{itemize}
        \item \lstinline{std::weak_ptr} is similar to \lstinline{std::shared_ptr},
            except it does \textbf{not} alter the reference count of the instance \pause
        \item A \lstinline{weak_ptr} must be \textbf{locked} in order to access the instance to which
            it points \pause
    \end{itemize}
    \begin{lstlisting}
std::shared_ptr<MyClass> lockedInstance
    = weakPtr.lock();
    \end{lstlisting} \pause
    \begin{itemize}
        \item Locking creates a temporary \lstinline{shared_ptr} to the instance ---
            this ensures that the reference count stays above 0 while the instance is being used
    \end{itemize}
\end{frame}

\begin{frame}[fragile]{Smart pointers --- Summary}
    \begin{itemize}
        \item \lstinline{std::shared_ptr} uses \textbf{reference counting} to manage the lifetime
            of objects \pause
        \item \lstinline{std::weak_ptr} is useful for breaking \textbf{reference cycles}
            when using \lstinline{shared_ptr}s for composition \pause
        \item \lstinline{std::unique_ptr} can be used when you want to enforce that there is
            \textbf{only one} pointer to a particular instance \pause
        \item NB: all smart pointer types allow \textbf{polymorphism}, e.g.
    \end{itemize}
    \begin{lstlisting}
class Bird {};
class Duck : public Bird {};

std::shared_ptr<Bird> bird = std::make_unique<Duck>();
    \end{lstlisting}
\end{frame}

\begin{frame}[fragile]{What about ``dumb'' pointers?}
    \begin{itemize}
        \item In modern C++, it is \textbf{almost always preferable} to use smart pointers
            instead of normal C++ pointers \pause
        \item I.e.\ use \lstinline{std::shared_ptr<T>} instead of \lstinline{T*},
            and \lstinline{std::make_shared} instead of \lstinline{new},
            and then you never have to worry about calling \lstinline{delete} \pause
        \item So why did we bother teaching you about normal C++ pointers?! \pause
            \begin{itemize}
                \item Because C++ code written before 2011 (and much code written after)
                    uses them extensively \pause
                \item Because sometimes you need to use them,
                    e.g.\ when working with SDL or other APIs / libraries \pause
                \item Because once you can write a nontrivial program that uses pointers and doesn't
                    crash or leak memory, you can call yourself a Real Programmer \smiley
            \end{itemize}
    \end{itemize}
\end{frame}

